\section{Grade two and grade four}\label{sec:grade-four}

\subsection{The grade-two zero fiber}

Let $g=\sum_{a=1}^3 e_a|e_a$.  Under the standard $SO(3)$ decomposition,
\begin{equation}\label{eq:b2-decomp}
  B_2=V\otimes V
  \cong \R g\oplus\Lambda^2V\oplus S^2_0V.
\end{equation}
Via the Euclidean metric, we identify $S^2_0V$ with the symmetric trace-free endomorphisms of $V$.
Let
\[
  P=\R g\oplus\Lambda^2V\cong V_0\oplus V_1,
  \qquad
  R=S^2_0V\cong V_2.
\]
The compression $m_2$ detects the trace and antisymmetric components and kills the symmetric trace-free component.  Hence
\begin{equation}\label{eq:grade-two-kernel}
  \ker m_2=R=S^2_0V.
\end{equation}
For $S\in R$, a direct quaternionic calculation gives
\begin{equation}\label{eq:grade-two-response}
  A_2(S)(d)=-2S(d).
\end{equation}
Thus every nonzero $S\in R$ has $\depth(S)=1$.  Equivalently, the grade-two residual sector is zero-valued but contextually active.

For $S,T\in R\setminus\{0\}$, Theorem~\ref{thm:additive-depth} gives
\begin{equation}\label{eq:pure-residual-depth}
  \depth(S\concat T)=2.
\end{equation}
The factorwise two-probe response can be written
\begin{equation}\label{eq:D13}
  D_{13}(S\concat T)(d_1,d_2)=4T(d_2)S(d_1),
\end{equation}
and is nonzero for suitable probes.

\subsection{Algebraic coherence delay}

The grade-four space splits by the two grade-two factors:
\begin{equation}\label{eq:b4-sector}
  B_4=(P\otimes P)\oplus(P\otimes R)\oplus(R\otimes P)\oplus(R\otimes R).
\end{equation}
\begin{proposition}[Grade-four depth table]\label{prop:grade-four-table}
The combined response ranks are
\begin{equation}\label{eq:b4-ranks}
\begin{array}{c|cccc}
 d&0&1&2&3\\
 \hline
 \operatorname{rank}D_4^{\le d}&4&32&72&81\\
 \dim F_d^{(4)}&77&49&9&0.
\end{array}
\end{equation}
The sectorwise invisible dimensions are
\begin{equation}\label{eq:b4-sector-table}
\begin{array}{c|rrrr|r}
 &P\otimes P&P\otimes R&R\otimes P&R\otimes R&\text{total}\\
\hline
F_0^{(4)}&12&20&20&25&77\\
F_1^{(4)}&0&12&12&25&49\\
F_2^{(4)}&0&0&0&9&9\\
F_3^{(4)}&0&0&0&0&0.
\end{array}
\end{equation}
Moreover,
\begin{equation}\label{eq:spin-four-survivor}
  F_2^{(4)}=F_2^{(4)}\cap(R\otimes R)=S^4_0V\cong V_4,
  \qquad \dim S^4_0V=9.
\end{equation}
\end{proposition}

\begin{proof}
The rank and sector-kernel tables are obtained from the exact integer matrices of all partial responses through each depth; no floating-point rank calculation is used.  The computation is reproduced by \texttt{scripts/verify\_depth4.py}.

For the final identification, $S^4_0V\cong V_4$ is the Cartan-product summand of
$R\otimes R=V_2\otimes V_2$.  By Lemma~\ref{lem:partial-equivariance}, every partial response is equivariant.  Moreover,
\[
  \Hom(V^{\otimes r},\Hq)\cong V^{\otimes r}\otimes(V_0\oplus V_1)
\]
contains no spin larger than $r+1$.  Thus depth zero has target spin at most one, depth one at most two, and depth two at most three.  Hence every such response vanishes on $V_4$, so $S^4_0V\subseteq F_2^{(4)}$.  The rank table gives $\dim F_2^{(4)}=9=\dim V_4$, proving equality.  Finally $F_3^{(4)}=0$ by all-gap reconstruction.
\end{proof}

By the factorization law, the depth-two responses with gap sets $\{1,2\}$ and $\{2,3\}$ vanish identically on $R\otimes R$ because one grade-two factor remains unprobed and is killed by $m_2$.  Hence the only potentially nonzero depth-two channel on this sector is \eqref{eq:D13}.  Since $F_3^{(4)}=0$, every nonzero vector in the subspace \eqref{eq:spin-four-survivor} has depth three.

\begin{definition}[Algebraic coherence delay]
A non-decomposable linear combination has algebraic coherence delay when it is detected strictly later than every nonzero decomposable tensor from which it is assembled.
\end{definition}

Equation \eqref{eq:spin-four-survivor} supplies such a delay: no nonzero pure tensor $S\concat T$ belongs to $F_2^{(4)}$ by \eqref{eq:pure-residual-depth}, but suitable linear combinations of these depth-two tensors do.  Their entire response profile through depth two cancels, and the combined state first appears at depth three.  This is an algebraic statement, not a claim of quantum coherence.

\begin{corollary}[Equivariant depth layers]\label{cor:depth-layers}
Let
\[
  L_d^{(4)}=F_{d-1}^{(4)}/F_d^{(4)},
  \qquad F_{-1}^{(4)}=B_4.
\]
Sector by sector, the first-detection layers are
\begin{equation}\label{eq:sector-layers}
\begin{array}{c|cccc}
 &L_0&L_1&L_2&L_3\\
\hline
P\otimes P&V_0\oplus V_1&V_0\oplus2V_1\oplus V_2&0&0\\
P\otimes R&0&V_1\oplus V_2&V_2\oplus V_3&0\\
R\otimes P&0&V_1\oplus V_2&V_2\oplus V_3&0\\
R\otimes R&0&0&V_0\oplus V_1\oplus V_2\oplus V_3&V_4.
\end{array}
\end{equation}
Consequently,
\begin{equation}\label{eq:layers}
\begin{aligned}
 L_0^{(4)}&\cong V_0\oplus V_1,\\
 L_1^{(4)}&\cong V_0\oplus4V_1\oplus3V_2,\\
 L_2^{(4)}&\cong V_0\oplus V_1\oplus3V_2\oplus3V_3,\\
 L_3^{(4)}&\cong V_4.
\end{aligned}
\end{equation}
\end{corollary}

\begin{proof}
We have $P\cong V_0\oplus V_1$ and $R\cong V_2$.  The restriction $m_2|_P:P\to\Hq$ is an equivariant isomorphism.  Hence the depth-zero quotient of $P\otimes P$ is $V_0\oplus V_1$; subtracting it from
$P\otimes P\cong2V_0\oplus3V_1\oplus V_2$ gives its depth-one row.

Each mixed sector decomposes as
$V_1\oplus2V_2\oplus V_3$.  Its only nonzero depth-one response has target
$\Hom(V,\Hq)$, whose maximum spin is two, so $V_3$ is in the depth-one kernel.  The twelve-dimensional sector kernel in Proposition~\ref{prop:grade-four-table} must therefore be $V_2\oplus V_3$, yielding the two mixed rows.  Finally,
$R\otimes R\cong V_0\oplus V_1\oplus V_2\oplus V_3\oplus V_4$, and Proposition~\ref{prop:grade-four-table} identifies its depth-three survivor as $V_4$.  This proves \eqref{eq:sector-layers}; summing the rows gives \eqref{eq:layers}.
\end{proof}

The layer dimensions are $4,28,40,9$.  In particular, copies of $V_0,V_1,V_2$ occur at more than one detection depth.  Representation type alone therefore does not determine visibility depth; factor origin also matters.
